% from ega4-8-fr.tex
\oldpage[IV-3]{43}
$f^{-1}(y)=\operatorname{Spec}(\mathcal A_y/\mathfrak m_y\mathcal A_y)$
(II, 1.5.5) est tel que l'homomorphisme correspondant
$k(y)=\mathcal O_y/\mathfrak m_y\to
\mathcal A_y/\mathfrak m_y\mathcal A_y$ soit bijectif~; comme $\mathcal A_y$
est un $\mathcal O_y$-module de type fini, le lemme de Nakayama montre que
$\mathcal O_y\to\mathcal A_y$ est surjectif, ce qui achève la démonstration.

\begin{remark}[8.11.5.1]
\label{IV.8.11.5.1-fr}
On notera que le raisonnement précédent prouve que si $f$ est un
\emph{monomorphisme}, alors, pour tout $y\in Y$, $f^{-1}(y)$ est vide ou
$k(y)$-isomorphe à $\operatorname{Spec}(k(y))$.
\end{remark}

\begin{proposition}[8.11.6]
\label{IV.8.11.6-fr}
Si un morphisme $f:X\to Y$ de présentation finie est un homéomorphisme
universel, il est fini, surjectif et radiciel (la réciproque étant vraie
d'après (2.4.5, (i))).
\end{proposition}

\begin{proof}
En effet, $f$ étant de type fini, universellement fermé, et séparé en vertu de
(2.4.4), est propre par définition (II, 5.4.1). Comme il est évidemment
quasi-fini (II, 6.2.3), il est fini d'après (8.11.1). On sait par ailleurs
qu'il est radiciel (2.4.4), et évidemment surjectif.
\end{proof}

\subsection{Nouvelle démonstration et généralisation du «~Main Theorem~» de Zariski}
\label{subsection:IV.8.12-fr}
\label{IV.8.12-fr}

\begin{lemma}[8.12.1]
\label{IV.8.12.1-fr}
Soient $f:X\to Y$ un morphisme quasi-compact et quasi-séparé, $\mathcal C$ une
$\mathcal O_Y$-Algèbre quasi-cohérente, $Z=\operatorname{Spec}(\mathcal C)$,
qui est un $Y$-préschéma affine au-dessus de $Y$. Soient $g:X\to Z$ un
$Y$-morphisme, $\varphi=\mathcal A(g):\mathcal C\to f_*(\mathcal O_X)$ le
$\mathcal O_Y$-homomorphisme correspondant de $\mathcal O_Y$-Algèbres
(II, 1.2.7). Supposons que $g$ soit une immersion. Alors, pour que l'image
fermée de $X$ par $g$ (I, 9.5.3) soit égale à $Z$, il faut et il suffit que
$\varphi$ soit injectif~; $g$ est alors une immersion ouverte.
\end{lemma}

\begin{proof}
L'hypothèse entraîne que $f_*(\mathcal O_X)$ est une $\mathcal O_Y$-Algèbre
quasi-cohérente (1.7.5)~; en outre, comme le morphisme canonique $h:Z\to Y$
est affine, donc quasi-compact et séparé, $g$ est un
morphisme quasi-compact et quasi-séparé (1.2.2 et 1.1.2), donc $g_*(\mathcal O_X)$ est une
$\mathcal O_Z$-Algèbre quasi-cohérente (1.7.5). Cela étant, dire que l'image
fermée de $X$ par $g$ est égale à $Z$ signifie (I, 9.5.2) que
l'homomorphisme canonique $\theta:\mathcal O_Z\to g_*(\mathcal O_X)$ est
injectif. Mais on a $h_*(\mathcal O_Z)=\mathcal C$ par définition de $Z$
(II, 1.3.1), et $h_*(g_*(\mathcal O_X))=f_*(\mathcal O_X)$. Comme $Z$ est
affine au-dessus de $Y$, il revient au même de dire que l'homomorphisme
$\theta:\mathcal O_Z\to g_*(\mathcal O_X)$ est injectif ou que
l'homomorphisme correspondant
$\varphi=h_*(\theta):\mathcal C\to f_*(\mathcal O_X)$ est injectif
(I, 1.3.9). Le fait que $g$ est alors une immersion ouverte résulte de
(I, 9.5.10) et de l'hypothèse que $g$ est une immersion.
\end{proof}

\begin{lemma}[8.12.2]
\label{IV.8.12.2-fr}
Soient $Y$ un préschéma quasi-compact et quasi-séparé, $f:X\to Y$ un
morphisme quasi-séparé de type fini, $\mathcal C$ une $\mathcal O_Y$-Algèbre
quasi-cohérente, $Z=\operatorname{Spec}(\mathcal C)$. Soient $g:X\to Z$ un
$Y$-morphisme, $\varphi:\mathcal C\to f_*(\mathcal O_X)$ le
$\mathcal O_Y$-homomorphisme correspondant de $\mathcal O_Y$-Algèbres. Soit
$(\mathcal C_\lambda)$ la famille filtrante croissante des sous-$\mathcal O_Y$-Algèbres
quasi-cohérentes de type fini de $\mathcal C$ (dont $\mathcal C$ est la
réunion ((I, 9.6.6) et (1.7.9)))~; posons
$Z_\lambda=\operatorname{Spec}(\mathcal C_\lambda)$ et soit $g_\lambda$ le
morphisme composé $X\xrightarrow{g}Z\to Z_\lambda$. Alors les conditions
suivantes sont équivalentes~:
\begin{enumerate}
\item[\emph{a)}] $g$ est une immersion.
\item[\emph{b)}] Il existe $\lambda$ tel que $g_\lambda$ soit une immersion.
\end{enumerate}
